\chapter{Macierze odwzorowań liniowych i zmiana bazy}
\label{ch:lin4-macierze-zmiana-bazy}
\index{macierz@Macierz}
\index{macierz odwzorowania@Macierz odwzorowania liniowego}
\index{zmiana bazy@Zmiana bazy}

\noindent\textbf{Cel rozdziału.}
W Rozdziale~\ref{ch:lin-przestrzenie-wektorowe}
zdefiniowaliśmy współrzędne wektorów,
a w Rozdziale~\ref{ch:lin3-odwzorowania-liniowe}
zbadaliśmy odwzorowania liniowe, ich jądra i obrazy.
Macierz pozwala zapisać każde odwzorowanie liniowe
między przestrzeniami skończenie wymiarowymi
za pomocą skończonej tablicy skalarów.
Pokażemy, dlaczego kolumny macierzy są obrazami
wektorów bazowych, skąd bierze się wzór na mnożenie
macierzy i jak zmienia się zapis operatora
przy zmianie baz.

W całym rozdziale \(K\) jest ciałem,
\(0_K\neq1_K\), a wszystkie przestrzenie wektorowe
rozpatrujemy nad \(K\). Jeśli w danym twierdzeniu
występują dwie bazy, traktujemy je jako
\emph{uporządkowane} rodziny wektorów.
Przyjmujemy konwencję, że współrzędne wektora
zapisujemy w kolumnie. Złożenie \(S\circ T\)
oznacza najpierw zastosowanie \(T\), a potem \(S\).
Wyniki tego rozdziału dotyczą wymiarów skończonych;
ich dowody nie wymagają aksjomatu wyboru.

\section{Macierze jako prostokątne tablice skalarów}
\label{sec:lin4-definicja-macierzy}

\begin{definicja}[Macierz i jej rozmiar]
\label{def:lin4-macierz}
Niech \(m,n\in\mathbb N_0\).
\emph{Macierzą} o \(m\) wierszach i \(n\) kolumnach
nad \(K\) nazywamy rodzinę
\(A=(a_{ij})_{1\leq i\leq m,\ 1\leq j\leq n}\)
skalarów \(a_{ij}\in K\).
Zbiór wszystkich takich macierzy oznaczamy
przez \(\operatorname{Mat}_{m\times n}(K)\).
Dwie macierze są równe wtedy i tylko wtedy,
gdy mają ten sam rozmiar oraz odpowiadające
sobie wyrazy są równe.
Macierze o \(m=n\) nazywamy \emph{kwadratowymi}.
\end{definicja}

Jeśli \(m=0\) albo \(n=0\), zbiór par
indeksów jest pusty, więc
\(\operatorname{Mat}_{m\times n}(K)\) zawiera
dokładnie jedną macierz. Nie mylimy jej
z macierzą o dodatnim rozmiarze
i wszystkich wyrazach zerowych.
Takie przypadki będą przydatne
przy odwzorowaniach przestrzeni zerowych.

\begin{definicja}[Dodawanie i mnożenie przez skalar]
\label{def:lin4-dzialania-punktowe}
Niech \(m,n\in\mathbb N_0\).
Dla \(A=(a_{ij}),B=(b_{ij})\)
w \(\operatorname{Mat}_{m\times n}(K)\)
oraz \(c\in K\) definiujemy
\[
 A+B=(a_{ij}+b_{ij}),\qquad cA=(ca_{ij}).
\]
Macierz, której wszystkie wyrazy są równe \(0_K\),
oznaczamy przez \(0_{m\times n}\).
Dla \(m=n\) piszemy również \(0_n\),
gdy kontekst nie powoduje niejednoznaczności.
\end{definicja}

\begin{stwierdzenie}[Przestrzeń wektorowa macierzy]
\label{prop:lin4-przestrzen-macierzy}
Niech \(m,n\in\mathbb N_0\).
Zbiór \(\operatorname{Mat}_{m\times n}(K)\)
z działaniami z Definicji~\ref{def:lin4-dzialania-punktowe}
jest przestrzenią wektorową nad \(K\)
o wymiarze \(mn\).
Dla \(m=0\) lub \(n=0\) jest przestrzenią zerową.
\end{stwierdzenie}
\begin{proof}
Dodawanie i mnożenie przez skalar wykonujemy
osobno na każdym wyrazie.
Na przykład dla \(A,B,C\) tego samego rozmiaru
i \(c,d\in K\) wyraz \((i,j)\) macierzy
\(c(A+B)+dC\) wynosi
\(c(a_{ij}+b_{ij})+dc_{ij}\).
Równości wynikające z aksjomatów przestrzeni
wektorowej nad \(K\) zachodzą na każdym wyrazie,
a zatem dla całych macierzy.
Wektorem zerowym jest \(0_{m\times n}\).

Dla \(1\leq i\leq m\), \(1\leq j\leq n\)
niech \(E_{ij}\) będzie \emph{macierzą bazową}
mającą \(1_K\) w miejscu \((i,j)\) i \(0_K\)
we wszystkich innych miejscach.
Każda macierz \(A\) ma rozkład
\[
 A=\sum_{i=1}^{m}\sum_{j=1}^{n}a_{ij}E_{ij}.
\]
Jeżeli \(\sum_{i,j}\lambda_{ij}E_{ij}=0_{m\times n}\),
to porównanie wyrazów \((i,j)\) daje
\(\lambda_{ij}=0_K\) dla każdej pary indeksów.
Rodzina \((E_{ij})\) jest więc bazą
o \(mn\) elementach.

Gdy \(m=0\) lub \(n=0\), nie ma żadnej pary indeksów,
a jedyna macierz jest zarazem wektorem zerowym.
Pusta rodzina stanowi jej bazę; wymiar wynosi
\(0=mn\).
\end{proof}

\begin{przyklad}[Rozkład macierzy na macierze bazowe]
\label{prz:lin4-rozklad-na-eij}
Nad \(\mathbb R\) rozważmy
\[
 A=\begin{pmatrix}2&-1&0\\3&4&5\end{pmatrix}.
\]
W przestrzeni \(\operatorname{Mat}_{2\times3}(\mathbb R)\)
mamy
\(A=2E_{11}-E_{12}+3E_{21}+4E_{22}+5E_{23}\).
Macierz \(E_{13}\) występuje ze współczynnikiem
zerowym. Współczynniki w tym rozkładzie są
jednoznaczne, bo macierze \(E_{ij}\) stanowią bazę.
\end{przyklad}

\section{Iloczyn macierzy i kolumny}
\label{sec:lin4-macierz-kolumna}

\begin{definicja}[Mnożenie macierzy przez kolumnę]
\label{def:lin4-iloczyn-kolumna}
Niech \(m,n\in\mathbb N_0\),
\(A=(a_{ij})\in\operatorname{Mat}_{m\times n}(K)\)
oraz \(x=(x_1,\ldots,x_n)^{\mathsf T}\in K^n\).
Przez \(Ax\in K^m\) rozumiemy kolumnę
o współrzędnych
\[
 (Ax)_i=\sum_{j=1}^{n}a_{ij}x_j,
 \qquad 1\leq i\leq m.
\]
Suma przy \(n=0\) jest pusta i wynosi \(0_K\).
\end{definicja}

\begin{stwierdzenie}[Liniowość i kolumny macierzy]
\label{prop:lin4-kolumny}
Niech \(m,n\in\mathbb N_0\) i
\(A\in\operatorname{Mat}_{m\times n}(K)\).
Odwzorowanie
\[
 L_A\colon K^n\longrightarrow K^m,\qquad x\longmapsto Ax
\]
jest liniowe. Jeżeli \(e_1,\ldots,e_n\)
oznaczają wektory bazy standardowej \(K^n\),
to \(L_A(e_j)\) jest \(j\)-tą kolumną \(A\).
W szczególności każde \(x\in K^n\) spełnia
\[
 Ax=\sum_{j=1}^{n}x_jL_A(e_j).
\]
\end{stwierdzenie}
\begin{proof}
Dla \(x,y\in K^n\), \(c,d\in K\)
oraz każdego \(1\leq i\leq m\)
z rozdzielności w ciele otrzymujemy
\[
 \begin{aligned}
 (A(cx+dy))_i
 &=\sum_{j=1}^{n}a_{ij}(cx_j+dy_j)\\
 &=c\sum_{j=1}^{n}a_{ij}x_j+
   d\sum_{j=1}^{n}a_{ij}y_j\\
 &=(cAx+dAy)_i.
 \end{aligned}
\]
Wszystkie współrzędne są równe, więc
\(L_A(cx+dy)=cL_A(x)+dL_A(y)\).

Wektor \(e_j\) ma współrzędną \(1_K\)
w miejscu \(j\) i zera w pozostałych,
stąd \((Ae_j)_i=a_{ij}\).
Ostatnia równość wynika z rozkładu
\(x=\sum_j x_je_j\) i liniowości \(L_A\).
Jeżeli \(n=0\), suma jest pusta;
jedyny wektor \(K^0\) przechodzi na zero \(K^m\),
zgodnie ze wzorem.
\end{proof}

Od tej chwili przez \(E_n=(e_1,\ldots,e_n)\)
oznaczamy \emph{uporządkowaną bazę standardową}
przestrzeni \(K^n\); przy \(n=0\) jest ona pusta.
Nie należy mylić rodziny \(E_n\)
z macierzami bazowymi \(E_{ij}\).

\begin{przyklad}[Obliczenie iloczynu macierzy i kolumny]
\label{prz:lin4-iloczyn-kolumny}
Dla macierzy
\[
 A=\begin{pmatrix}1&1&0\\0&1&1\end{pmatrix},
 \qquad
 x=\begin{pmatrix}2\\-1\\3\end{pmatrix}
\]
nad \(\mathbb R\) otrzymujemy
\[
 Ax=\begin{pmatrix}
 1\cdot2+1\cdot(-1)+0\cdot3\\
 0\cdot2+1\cdot(-1)+1\cdot3
 \end{pmatrix}
 =\begin{pmatrix}1\\2\end{pmatrix}.
\]
Iloczyn jest określony, bo liczba kolumn
macierzy \(A\) jest równa liczbie
współrzędnych kolumny \(x\).
\end{przyklad}

\section{Macierz odwzorowania liniowego w wybranych bazach}
\label{sec:lin4-macierz-odwzorowania}

\begin{definicja}[Macierz odwzorowania w bazach]
\label{def:lin4-macierz-odwzorowania}
Niech \(V,W\) będą przestrzeniami wektorowymi
nad \(K\) o wymiarach odpowiednio \(n,m\in\mathbb N_0\).
Niech \(B=(b_1,\ldots,b_n)\) będzie uporządkowaną
bazą \(V\), \(C=(c_1,\ldots,c_m)\) uporządkowaną
bazą \(W\), a \(T\colon V\to W\) odwzorowaniem liniowym.
\emph{Macierzą odwzorowania \(T\) w bazach \(B,C\)}
nazywamy macierz
\[
 [T]_{C\leftarrow B}=(a_{ij})
 \in\operatorname{Mat}_{m\times n}(K),
 \qquad
 T(b_j)=\sum_{i=1}^{m}a_{ij}c_i.
\]
\(j\)-ta kolumna składa się więc ze współrzędnych
\([T(b_j)]_C\). Strzałka w indeksie
\(C\leftarrow B\) przypomina, że współrzędne
wejściowe odnoszą się do \(B\),
a wyjściowe do \(C\).
\end{definicja}

\begin{twierdzenie}[Wzór współrzędnych i jednoznaczność macierzy]
\label{thm:lin4-wspolrzedne-macierz}
Niech \(V,W\) będą przestrzeniami wektorowymi
nad \(K\), o wymiarach odpowiednio
\(n,m\in\mathbb N_0\). Niech
\(B=(b_1,\ldots,b_n)\) i \(C=(c_1,\ldots,c_m)\)
będą ich uporządkowanymi bazami,
a \(T\colon V\to W\) odwzorowaniem liniowym.
Wtedy dla każdego \(v\in V\)
\[
 [T(v)]_C=[T]_{C\leftarrow B}\,[v]_B.
\]
Ponadto \([T]_{C\leftarrow B}\) jest jedyną
macierzą \(A\in\operatorname{Mat}_{m\times n}(K)\),
dla której \([T(v)]_C=A[v]_B\)
dla wszystkich \(v\in V\).
Odwrotnie, każda macierz
\(A\in\operatorname{Mat}_{m\times n}(K)\)
jest macierzą dokładnie jednego odwzorowania
liniowego \(T\colon V\to W\) w tych bazach.
\end{twierdzenie}
\begin{proof}
Niech \([v]_B=(x_1,\ldots,x_n)^{\mathsf T}\).
Jednoznaczność współrzędnych z
Twierdzenia~\ref{thm:lin-wspolrzedne}
daje \(v=\sum_{j=1}^{n}x_jb_j\).
Z liniowości
\[
 \begin{aligned}
 T(v)&=\sum_{j=1}^{n}x_jT(b_j)\\
 &=\sum_{j=1}^{n}x_j
     \sum_{i=1}^{m}a_{ij}c_i
 =\sum_{i=1}^{m}
    \left(\sum_{j=1}^{n}a_{ij}x_j\right)c_i.
 \end{aligned}
\]
Przestawiamy tu wyłącznie sumy skończone.
Współczynnik przy \(c_i\) jest
\(i\)-tą współrzędną iloczynu macierzy
\([T]_{C\leftarrow B}[v]_B\).
Jednoznaczność zapisu w bazie \(C\)
dowodzi pierwszego wzoru.

Jeśli inna macierz \(A\) spełnia go
dla wszystkich \(v\), to podstawiając
\(v=b_j\) otrzymujemy
\[
 A e_j=[T(b_j)]_C,
\]
gdyż \([b_j]_B=e_j\).
\(j\)-ta kolumna \(A\) jest więc wyznaczona
jednoznacznie i musi być równa kolumnie
\([T]_{C\leftarrow B}\). Stąd równość macierzy.

Dla dowolnego \(A\) definiujemy
\[
 T(v):=[\,\cdot\,]_C^{-1}
         \bigl(A[v]_B\bigr),
 \qquad v\in V,
\]
gdzie \([\,\cdot\,]_C^{-1}\colon K^m\to W\)
jest odwrotnością izomorfizmu współrzędnych.
Składamy trzy odwzorowania liniowe,
zatem \(T\) jest liniowe.
Z konstrukcji \([T(v)]_C=A[v]_B\),
a wykazana jednoznaczność kończy dowód.

W przypadkach \(n=0\) i \(m=0\)
wszystkie sumy są rozumiane jako puste.
Istnieje wówczas jedna odpowiednia macierz
i jedno odwzorowanie liniowe, więc
argument pozostaje prawdziwy.
\end{proof}

\begin{figure}[!htbp]
\centering
\begin{tikzpicture}[>=Latex,every node/.style={font=\small}]
 \node (V) at (0,2.1) {\(V\)};
 \node (W) at (5.2,2.1) {\(W\)};
 \node (Kn) at (0,0) {\(K^n\)};
 \node (Km) at (5.2,0) {\(K^m\)};
 \draw[->,thick] (V)--node[above]{\(T\)}(W);
 \draw[->,thick] (V)--node[left]{\([\,\cdot\,]_B\)}(Kn);
 \draw[->,thick] (W)--node[right]{\([\,\cdot\,]_C\)}(Km);
 \draw[->,thick] (Kn)--node[below]{\(L_A\)}(Km);
\end{tikzpicture}
\caption{Diagram współrzędnych dla
\(A=[T]_{C\leftarrow B}\).
Obie drogi z \(V\) do \(K^m\) dają tę samą kolumnę:
\([T(v)]_C=A[v]_B\).}
\label{fig:lin4-wspolrzedne}
\end{figure}

\begin{przyklad}[Macierz z obrazu wektorów bazowych]
\label{prz:lin4-macierz-z-obrazow}
Niech \(T\colon\mathbb R^3\to\mathbb R^2\)
będzie dane przez \(T(x,y,z)=(x+y,y+z)\),
tak jak w Przykładzie~\ref{prz:lin3-rzad-r3}.
Dla baz standardowych \(E_3\) i \(E_2\) mamy
\[
 T(e_1)=\begin{pmatrix}1\\0\end{pmatrix},
 \quad T(e_2)=\begin{pmatrix}1\\1\end{pmatrix},
 \quad T(e_3)=\begin{pmatrix}0\\1\end{pmatrix}.
\]
Ustawiając te kolumny w podanej kolejności,
otrzymujemy
\[
 [T]_{E_2\leftarrow E_3}
 =\begin{pmatrix}1&1&0\\0&1&1\end{pmatrix}.
\]
Obliczenie na kolumnie \((2,-1,3)^{\mathsf T}\)
z Przykładu~\ref{prz:lin4-iloczyn-kolumny}
daje \((1,2)^{\mathsf T}\), co zgadza się
z \(T(2,-1,3)=(1,2)\).
\end{przyklad}

\begin{twierdzenie}[Izomorfizm przestrzeni odwzorowań i macierzy]
\label{thm:lin4-hom-mat}
Niech \(V,W\) będą przestrzeniami wektorowymi
nad \(K\) o skończonych wymiarach odpowiednio
\(n,m\in\mathbb N_0\), a \(B,C\)
ich ustalonymi uporządkowanymi bazami.
Wtedy przyporządkowanie
\[
 \mathcal M_{C\leftarrow B}\colon
 \operatorname{Hom}_K(V,W)
 \longrightarrow\operatorname{Mat}_{m\times n}(K),
 \qquad T\longmapsto[T]_{C\leftarrow B}
\]
jest izomorfizmem przestrzeni wektorowych.
W szczególności
\(\dim_K\operatorname{Hom}_K(V,W)=mn\).
\end{twierdzenie}
\begin{proof}
Dla \(S,T\colon V\to W\) liniowych
oraz \(a,b\in K\), z równości
\((aS+bT)(b_j)=aS(b_j)+bT(b_j)\)
wynika, że każda kolumna macierzy
\([aS+bT]_{C\leftarrow B}\)
jest odpowiednią kombinacją kolumn
macierzy \(S\) i \(T\).
Zatem
\[
 [aS+bT]_{C\leftarrow B}
 =a[S]_{C\leftarrow B}
  +b[T]_{C\leftarrow B},
\]
co dowodzi liniowości \(\mathcal M_{C\leftarrow B}\).
Poprzednie twierdzenie zapewnia, że
każda macierz pochodzi od dokładnie jednego \(T\).
Odwzorowanie jest więc bijektywne.
Jako liniowa bijekcja jest izomorfizmem
na mocy Twierdzenia~\ref{thm:lin3-odwrotnosc}.
Równość wymiarów wynika z
Stwierdzenia~\ref{prop:lin4-przestrzen-macierzy}.
Dowód nie wymaga aksjomatu wyboru,
gdyż bazy \(B,C\) są już dane.
\end{proof}

\section{Mnożenie macierzy i składanie odwzorowań}
\label{sec:lin4-mnozenie}

\begin{definicja}[Iloczyn dwóch macierzy]
\label{def:lin4-iloczyn-macierzy}
Niech \(m,n,p\in\mathbb N_0\),
\(A=(a_{ij})\in\operatorname{Mat}_{m\times n}(K)\)
i \(B=(b_{jk})\in\operatorname{Mat}_{n\times p}(K)\).
\emph{Iloczynem} \(AB\) nazywamy macierz
w \(\operatorname{Mat}_{m\times p}(K)\)
o wyrazach
\[
 (AB)_{ik}=\sum_{j=1}^{n}a_{ij}b_{jk},
 \qquad 1\leq i\leq m,\quad 1\leq k\leq p.
\]
Wymiar wewnętrzny \(n\) musi być taki sam
dla obu macierzy. Iloczyn \(BA\) może
nie być określony.
\end{definicja}

\begin{twierdzenie}[Macierz złożenia]
\label{thm:lin4-macierz-zlozenia}
Niech \(U,V,W\) będą przestrzeniami
wektorowymi nad \(K\), o wymiarach
odpowiednio \(p,n,m\in\mathbb N_0\).
Niech \(D,B,C\) będą uporządkowanymi bazami
tych przestrzeni, a \(S\colon U\to V\),
\(T\colon V\to W\) odwzorowaniami liniowymi.
Wtedy
\[
 [T\circ S]_{C\leftarrow D}
 =[T]_{C\leftarrow B}\,[S]_{B\leftarrow D}.
\]
\end{twierdzenie}
\begin{proof}
Niech \(A=[T]_{C\leftarrow B}\),
\(M=[S]_{B\leftarrow D}\).
Dla dowolnego \(u\in U\)
dwukrotne zastosowanie
Twierdzenia~\ref{thm:lin4-wspolrzedne-macierz}
daje
\[
 [(T\circ S)(u)]_C
 =A[S(u)]_B
 =A(M[u]_D).
\]
Sprawdźmy, że \(A(Mx)=(AM)x\)
dla każdej kolumny \(x\in K^p\).
Jej \(i\)-ta współrzędna po lewej stronie
wynosi
\[
 \sum_{j=1}^{n}a_{ij}
   \left(\sum_{k=1}^{p}m_{jk}x_k\right)
 =\sum_{k=1}^{p}
   \left(\sum_{j=1}^{n}a_{ij}m_{jk}\right)x_k,
\]
co z definicji jest \(i\)-tą współrzędną
\((AM)x\). Przestawienie sum jest uprawnione,
ponieważ obie sumy są skończone.
Ostatecznie
\([(T\circ S)(u)]_C=(AM)[u]_D\).
Jednoznaczność macierzy odwzorowania
daje wzór z tezy.
\end{proof}

\begin{przyklad}[Mnożenie macierzy]
\label{prz:lin4-mnozenie}
Niech
\[
 A=\begin{pmatrix}1&2&0\\0&-1&3\end{pmatrix},
 \qquad
 B=\begin{pmatrix}1&0\\2&1\\0&-1\end{pmatrix}
\]
będą macierzami nad \(\mathbb R\),
odpowiednio o rozmiarach \(2\times3\)
i \(3\times2\). Wtedy
\[
 AB=
 \begin{pmatrix}
 1\cdot1+2\cdot2+0\cdot0&
 1\cdot0+2\cdot1+0\cdot(-1)\\
 0\cdot1+(-1)\cdot2+3\cdot0&
 0\cdot0+(-1)\cdot1+3\cdot(-1)
 \end{pmatrix}
 =\begin{pmatrix}5&2\\-2&-4\end{pmatrix}.
\]
Natomiast \(BA\) ma rozmiar \(3\times3\);
nie może więc być równe \(AB\)
nawet przed wykonaniem rachunków.
\end{przyklad}

\begin{definicja}[Macierz jednostkowa]
\label{def:lin4-jednostkowa}
Dla \(n\in\mathbb N_0\)
macierzą \emph{jednostkową} \(I_n\)
nazywamy macierz kwadratową
o wyrazach
\[
 (I_n)_{ij}=
 \begin{cases}
 1_K,&i=j,\\
 0_K,&i\neq j.
 \end{cases}
\]
Jej kolumny to wektory bazy standardowej \(K^n\).
Dla \(n=0\) jest to jedyna macierz \(0\times0\).
\end{definicja}

\begin{twierdzenie}[Łączność i rozdzielność iloczynu macierzy]
\label{thm:lin4-rachunek-macierzy}
Niech \(m,n,p,q\in\mathbb N_0\) oraz
\(A\in\operatorname{Mat}_{m\times n}(K)\),
\(B\in\operatorname{Mat}_{n\times p}(K)\),
\(C\in\operatorname{Mat}_{p\times q}(K)\).
Wtedy
\[
 (AB)C=A(BC).
\]
Ponadto dla dowolnych
\(A'\in\operatorname{Mat}_{m\times n}(K)\),
\(B'\in\operatorname{Mat}_{n\times p}(K)\)
i \(a\in K\) iloczyn jest rozdzielny
względem dodawania i zgodny z mnożeniem
przez skalar:
\[
 \begin{gathered}
 (A+A')B=AB+A'B,\qquad
 A(B+B')=AB+AB',\\
 (aA)B=a(AB)=A(aB).
 \end{gathered}
\]
Dla każdego \(D\in\operatorname{Mat}_{m\times n}(K)\)
zachodzi \(I_mD=D=DI_n\).
W przypadku macierzy kwadratowych
iloczyn na ogół nie jest przemienny.
\end{twierdzenie}
\begin{proof}
Wyraz \((i,\ell)\) macierzy \((AB)C\)
ma postać
\[
 \sum_{k=1}^{p}(AB)_{ik}c_{k\ell}
 =\sum_{k=1}^{p}\sum_{j=1}^{n}
    a_{ij}b_{jk}c_{k\ell}.
\]
W macierzy \(A(BC)\) ten sam wyraz
wynosi \(\sum_{j=1}^{n}\sum_{k=1}^{p}
a_{ij}b_{jk}c_{k\ell}\).
Równości są identyczne po przestawieniu
dwóch sum skończonych, co dowodzi łączności.

Dla przykładu wyraz \((i,k)\) iloczynu
\((A+A')B\) wynosi
\[
 \sum_{j=1}^{n}(a_{ij}+a'_{ij})b_{jk}
 =\sum_{j=1}^{n}a_{ij}b_{jk}
  +\sum_{j=1}^{n}a'_{ij}b_{jk}.
\]
Jest to wyraz macierzy \(AB+A'B\).
Pozostałe równości wynikają z tych samych aksjomatów ciała.
Dla macierzy \(A\) i \(B,B'\) odpowiednich rozmiarów
oraz skalara \(c\in K\) sprawdzamy ich współrzędne:
\[
 \begin{aligned}
 (A(B+B'))_{ik}&=\sum_j a_{ij}(b_{jk}+b'_{jk})
   =(AB)_{ik}+(AB')_{ik},\\
 ((cA)B)_{ik}&=\sum_j(ca_{ij})b_{jk}
   =c\sum_j a_{ij}b_{jk}=(c(AB))_{ik},\\
 (A(cB))_{ik}&=\sum_j a_{ij}(cb_{jk})
   =c\sum_j a_{ij}b_{jk}=(c(AB))_{ik}.
 \end{aligned}
\]
Każda suma jest skończona; wyciągnięcie skalara
przed sumę korzysta z rozdzielności w ciele \(K\).

Wyraz \((i,j)\) macierzy \(I_mD\)
jest równy
\(\sum_{k=1}^{m}(I_m)_{ik}d_{kj}=d_{ij}\),
ponieważ tylko składnik z \(k=i\)
może być niezerowy.
Podobnie \(DI_n=D\).
Jeśli któryś wymiar wynosi zero,
wszystkie odpowiednie równości dotyczą
jedynych macierzy rozmiaru z pustym
zbiorem indeksów albo sum pustych;
pozostają więc prawdziwe.

Nieprzemienność pokażemy już dla \(n=2\).
Niech
\[
 P=\begin{pmatrix}1&0\\0&0\end{pmatrix},
 \qquad
 Q=\begin{pmatrix}0&1\\1&0\end{pmatrix}.
\]
Wtedy
\[
 PQ=\begin{pmatrix}0&1\\0&0\end{pmatrix},
 \qquad
 QP=\begin{pmatrix}0&0\\1&0\end{pmatrix}.
\]
Macierze są różne nad każdym ciałem,
gdyż \(1_K\neq0_K\).
\end{proof}

\section{Jądro, obraz i rząd macierzy}
\label{sec:lin4-rzad-macierzy}

\begin{definicja}[Jądro, obraz i rząd macierzy]
\label{def:lin4-rzad-macierzy}
Niech \(m,n\in\mathbb N_0\) oraz
\(A\in\operatorname{Mat}_{m\times n}(K)\).
Definiujemy
\[
 \ker A:=\{x\in K^n:Ax=0_{K^m}\},\qquad
 \operatorname{im}A:=\{Ax:x\in K^n\},
\]
a \emph{rzędem macierzy} nazywamy liczbę
\[
 \operatorname{rk}A:=\dim_K(\operatorname{im}A).
\]
Są to odpowiednio jądro, obraz i rząd
odwzorowania \(L_A\colon K^n\to K^m\).
\end{definicja}

\begin{twierdzenie}[Kolumny, jądro i rząd]
\label{thm:lin4-rzad-kolumn}
Niech \(m,n\in\mathbb N_0\), a
\(A\in\operatorname{Mat}_{m\times n}(K)\)
ma kolumny \(a_1,\ldots,a_n\in K^m\).
Wtedy
\[
 \operatorname{im}A=\operatorname{span}\{a_1,\ldots,a_n\},
 \qquad
 \ker A=\left\{x\in K^n:
   \sum_{j=1}^{n}x_ja_j=0_{K^m}\right\}.
\]
W szczególności
\(\operatorname{rk}A\) jest wymiarem
przestrzeni rozpiętej przez kolumny,
a \(0\leq\operatorname{rk}A\leq\min\{m,n\}\).
Ponadto
\[
 n=\dim_K(\ker A)+\operatorname{rk}A.
\]
\end{twierdzenie}
\begin{proof}
Na mocy Stwierdzenia~\ref{prop:lin4-kolumny}
dla każdego \(x\in K^n\) zachodzi
\(Ax=\sum_{j=1}^{n}x_ja_j\).
Każdy element obrazu jest więc
kombinacją kolumn, a każda taka kombinacja
jest obrazem kolumny współczynników \(x\).
Dowodzi to pierwszej równości.
Druga jest definicją jądra po podstawieniu
tego samego wzoru.

Obraz jest podprzestrzenią \(K^m\),
zatem jego wymiar nie przekracza \(m\).
Jest też rozpięty przez \(n\) kolumn,
więc wymiar nie przekracza \(n\)
na mocy wyników o rodzinach rozpinających
z Rozdziału~\ref{ch:lin-przestrzenie-wektorowe}.
Twierdzenie~\ref{thm:lin3-rzad-nullity}
zastosowane do \(L_A\) daje
\(n=\dim_K\ker A+\operatorname{rk}A\).
Wszystkie wymiary są skończone.
\end{proof}

\begin{wniosek}[Kryteria iniektywności i surjektywności macierzy]
\label{cor:lin4-kryteria-kolumn}
Niech \(m,n\in\mathbb N_0\) oraz
\(A\in\operatorname{Mat}_{m\times n}(K)\).
Wtedy \(L_A\colon K^n\to K^m\) jest
iniektywne wtedy i tylko wtedy, gdy
kolumny \(A\) są liniowo niezależne.
Jest surjektywne wtedy i tylko wtedy,
gdy kolumny rozpinają \(K^m\).
\end{wniosek}
\begin{proof}
Z poprzedniego twierdzenia \(\ker A=\{0\}\)
dokładnie wtedy, gdy warunek
\(\sum_jx_ja_j=0\) wymusza \(x_j=0_K\)
dla wszystkich \(j\).
To definicja niezależności liniowej kolumn.
Surjektywność jest równoważna
\(\operatorname{im}A=K^m\),
czyli rozpinaniu przez kolumny.
\end{proof}

\begin{przyklad}[Rząd prostokątnej macierzy]
\label{prz:lin4-rzad-prostokatnej}
Dla
\[
 A=\begin{pmatrix}1&1&0\\0&1&1\end{pmatrix}
\]
nad \(\mathbb R\) pierwsza i trzecia kolumna
to \((1,0)^{\mathsf T}\) i \((0,1)^{\mathsf T}\).
Tworzą bazę \(\mathbb R^2\), zatem
\(\operatorname{im}A=\mathbb R^2\)
i \(\operatorname{rk}A=2\).
Układ \(Ax=0\) daje
\(x_1+x_2=0,\ x_2+x_3=0\), więc
\[
 \ker A=\operatorname{span}\{
 (-1,1,-1)^{\mathsf T}\}.
\]
Wymiar jądra wynosi \(1\), zgodnie z
\(3=1+2\).
\end{przyklad}

\section{Macierze odwracalne}
\label{sec:lin4-odwracalne}

\begin{definicja}[Macierz odwracalna]
\label{def:lin4-odwracalna}
Niech \(n\in\mathbb N_0\).
Macierz \(A\in\operatorname{Mat}_{n\times n}(K)\)
nazywamy \emph{odwracalną}, gdy istnieje
\(B\in\operatorname{Mat}_{n\times n}(K)\)
taka, że \(AB=I_n=BA\).
Macierz \(B\), jeśli istnieje, nazywamy
\emph{odwrotną} do \(A\) i oznaczamy \(A^{-1}\).
\end{definicja}

\begin{twierdzenie}[Kryteria odwracalności]
\label{thm:lin4-odwracalnosc-kryteria}
Niech \(n\in\mathbb N_0\) oraz
\(A\in\operatorname{Mat}_{n\times n}(K)\).
Następujące warunki są równoważne:
\begin{enumerate}[label=(\roman*)]
\item \(A\) jest odwracalna;
\item \(L_A\colon K^n\to K^n\) jest bijektywne;
\item kolumny \(A\) tworzą bazę \(K^n\);
\item \(\ker A=\{0_{K^n}\}\);
\item \(\operatorname{im}A=K^n\).
\end{enumerate}
Gdy zachodzą, macierz odwrotna jest jednoznaczna
i spełnia
\[
 [L_A^{-1}]_{E_n\leftarrow E_n}=A^{-1}.
\]
\end{twierdzenie}
\begin{proof}
Jeśli \(AB=I_n=BA\), to
z Twierdzenia~\ref{thm:lin4-macierz-zlozenia}
otrzymujemy
\(L_A\circ L_B=\operatorname{id}_{K^n}
 =L_B\circ L_A\).
Odwzorowanie \(L_A\) jest zatem bijektywne.

Jeśli \(L_A\) jest bijektywne,
jego odwrotność jest liniowa na mocy
Twierdzenia~\ref{thm:lin3-odwrotnosc}.
Niech \(B\) będzie macierzą \(L_A^{-1}\)
w bazie standardowej.
Wzór na macierz złożenia daje
\(AB=I_n=BA\). Stąd (i) i (ii)
są równoważne.

Kolumny \(A\) są obrazami bazy standardowej
przez \(L_A\). Na mocy
Stwierdzenia~\ref{prop:lin3-izomorfizm-baza}
bijektywność \(L_A\) implikuje, że kolumny
tworzą bazę. Odwrotnie, gdy kolumny
tworzą bazę, są niezależne i rozpinające,
więc zgodnie z Wnioskiem~\ref{cor:lin4-kryteria-kolumn}
odwzorowanie \(L_A\) jest zarówno iniektywne,
jak i surjektywne. Otrzymaliśmy (ii)
równoważne (iii).

Na mocy Wniosku~\ref{cor:lin3-rownymiary}
dla endomorfizmu skończenie wymiarowej
przestrzeni \(K^n\) iniektywność,
surjektywność i bijektywność są równoważne.
Twierdzenie~\ref{thm:lin3-iniekcja-surjekcja}
utożsamia iniektywność z (iv),
a surjektywność z (v).
Wszystkie warunki są zatem równoważne.

Jeżeli \(B,C\) są dwiema macierzami odwrotnymi
do \(A\), to z łączności mnożenia
\[
 B=B(AC)=(BA)C=I_nC=C.
\]
To dowodzi jednoznaczności.
Przy \(n=0\) wszystkie warunki zachodzą
dla jedynej macierzy \(0\times0\),
a wzory pozostają poprawne.
\end{proof}

\begin{przyklad}[Odwrotność macierzy o rozmiarze \(2\times2\)]
\label{prz:lin4-odwrotnosc-2}
Nad \(\mathbb R\) niech
\[
 A=\begin{pmatrix}2&1\\1&1\end{pmatrix},
 \qquad
 B=\begin{pmatrix}1&-1\\-1&2\end{pmatrix}.
\]
Mnożenie daje
\[
 AB=
 \begin{pmatrix}
 2\cdot1+1\cdot(-1)&2\cdot(-1)+1\cdot2\\
 1\cdot1+1\cdot(-1)&1\cdot(-1)+1\cdot2
 \end{pmatrix}
 =I_2.
\]
Dla iloczynu w odwrotnej kolejności obliczamy również
\[
 BA=\begin{pmatrix}
 1\cdot2+(-1)\cdot1&1\cdot1+(-1)\cdot1\\
 (-1)\cdot2+2\cdot1&(-1)\cdot1+2\cdot1
 \end{pmatrix}=I_2.
\]
Zatem \(A^{-1}=B\).
Nie skorzystaliśmy jeszcze z wyznaczników;
wystarczy sprawdzić oba iloczyny.
\end{przyklad}

\begin{stwierdzenie}[Iloczyn i odwrotność macierzy odwracalnych]
\label{prop:lin4-iloczyn-odwrotnych}
Niech \(n\in\mathbb N_0\), a
\(A,B\in\operatorname{Mat}_{n\times n}(K)\)
będą odwracalne. Wtedy \(AB\) jest odwracalna
oraz
\[
 (AB)^{-1}=B^{-1}A^{-1},\qquad
 (A^{-1})^{-1}=A.
\]
\end{stwierdzenie}
\begin{proof}
Korzystając z łączności mnożenia,
\[
 (AB)(B^{-1}A^{-1})
 =A(BB^{-1})A^{-1}=AA^{-1}=I_n.
\]
Podobnie
\((B^{-1}A^{-1})(AB)
 =B^{-1}(A^{-1}A)B=I_n\).
Jednoznaczność odwrotności daje pierwszy wzór.
Równości \(AA^{-1}=I_n=A^{-1}A\)
dowodzą, że \(A\) jest odwrotnością \(A^{-1}\),
czyli drugiego wzoru.
\end{proof}

\section{Macierze zmiany bazy}
\label{sec:lin4-zmiana-bazy}

\begin{definicja}[Macierz przejścia między bazami]
\label{def:lin4-przejscie}
Niech \(n\in\mathbb N_0\), a \(V\)
będzie \(n\)-wymiarową przestrzenią
wektorową nad \(K\). Niech
\(B=(b_1,\ldots,b_n)\) i
\(B'=(b'_1,\ldots,b'_n)\) jej uporządkowanymi bazami.
\emph{Macierzą przejścia od współrzędnych w \(B'\)
do współrzędnych w \(B\)} nazywamy
\[
 P_{B\leftarrow B'}
 :=\bigl([b'_1]_B\ \cdots\ [b'_n]_B\bigr)
 \in\operatorname{Mat}_{n\times n}(K).
\]
Jej \(j\)-ta kolumna zawiera współrzędne
\(j\)-tego wektora nowej bazy \(B'\)
w starej bazie \(B\).
\end{definicja}

\begin{twierdzenie}[Współrzędne przy zmianie bazy]
\label{thm:lin4-zmiana-wspolrzednych}
Niech \(V\) będzie przestrzenią wektorową
nad \(K\) o wymiarze \(n\in\mathbb N_0\),
a \(B=(b_1,\ldots,b_n)\) i
\(B'=(b'_1,\ldots,b'_n)\) jej
uporządkowanymi bazami.
Macierz \(P:=P_{B\leftarrow B'}\)
jest odwracalna i dla każdego \(v\in V\)
zachodzi
\[
 [v]_B=P[v]_{B'},\qquad
 [v]_{B'}=P^{-1}[v]_B.
\]
Ponadto
\[
 P_{B'\leftarrow B}=P_{B\leftarrow B'}^{-1}.
\]
\end{twierdzenie}
\begin{proof}
Rozważmy odwzorowanie identycznościowe
\(\operatorname{id}_V\colon V\to V\).
Z definicji jego macierz w bazach
wejściowej \(B'\) i wyjściowej \(B\)
ma w \(j\)-tej kolumnie \([b'_j]_B\),
a więc jest równa \(P\).
Twierdzenie~\ref{thm:lin4-wspolrzedne-macierz}
daje
\([v]_B=P[v]_{B'}\) dla każdego \(v\in V\).

Podobnie \(Q:=P_{B'\leftarrow B}\)
spełnia \([v]_{B'}=Q[v]_B\).
Po podstawieniu pierwszej równości
otrzymujemy
\([v]_{B'}=QP[v]_{B'}\).
Ponieważ odwzorowanie współrzędnych
\(V\to K^n\) jest surjektywne,
mamy \(QP x=x\) dla każdego \(x\in K^n\),
co daje \(QP=I_n\).
W odwrotnej kolejności podstawień
otrzymujemy \(PQ=I_n\).
Zatem \(P\) jest odwracalna i \(Q=P^{-1}\).
Wszystko pozostaje poprawne przy \(n=0\).
\end{proof}

\begin{przyklad}[Współrzędne w dwóch bazach]
\label{prz:lin4-zmiana-wspolrzednych}
Niech \(V=\mathbb R^2\),
\(B=((1,0),(0,1))\) oraz
\(B'=((1,1),(0,1))\).
Wtedy
\[
 P_{B\leftarrow B'}
 =\begin{pmatrix}1&0\\1&1\end{pmatrix},
 \qquad
 P_{B\leftarrow B'}^{-1}
 =\begin{pmatrix}1&0\\-1&1\end{pmatrix}.
\]
Dla wektora \(v=(2,3)\) mamy
\([v]_B=(2,3)^{\mathsf T}\),
a zatem
\[
 [v]_{B'}=
 \begin{pmatrix}1&0\\-1&1\end{pmatrix}
 \begin{pmatrix}2\\3\end{pmatrix}
 =\begin{pmatrix}2\\1\end{pmatrix}.
\]
Istotnie, \((2,3)=2(1,1)+(0,1)\).
Nie zmieniliśmy wektora \(v\), tylko
jego współrzędne.
\end{przyklad}

\begin{twierdzenie}[Zmiana baz dziedziny i przeciwdziedziny]
\label{thm:lin4-zmiana-macierzy}
Niech \(V,W\) będą przestrzeniami wektorowymi
nad \(K\), o wymiarach \(n,m\in\mathbb N_0\),
a \(T\colon V\to W\) odwzorowaniem liniowym.
Niech \(B,B'\) będą uporządkowanymi bazami \(V\),
a \(C,C'\) uporządkowanymi bazami \(W\).
Oznaczmy
\[
 A=[T]_{C\leftarrow B},\qquad
 P=P_{B\leftarrow B'},\qquad
 Q=P_{C\leftarrow C'}.
\]
Wówczas
\[
 [T]_{C'\leftarrow B'}=Q^{-1}AP.
\]
\end{twierdzenie}
\begin{proof}
Dla dowolnego \(v\in V\) mamy
\([v]_B=P[v]_{B'}\) na mocy
Twierdzenia~\ref{thm:lin4-zmiana-wspolrzednych}.
Ponadto
\([T(v)]_C=A[v]_B\).
W przypadku przestrzeni \(W\)
to samo twierdzenie daje
\([T(v)]_{C'}=Q^{-1}[T(v)]_C\).
Łącząc te równości, otrzymujemy
\[
 [T(v)]_{C'}=Q^{-1}AP[v]_{B'}.
\]
Macierz \(Q^{-1}AP\) ma rozmiar \(m\times n\).
Jednoznaczność macierzy odwzorowania
z Twierdzenia~\ref{thm:lin4-wspolrzedne-macierz}
daje tezę.
\end{proof}

\begin{wniosek}[Macierze podobne]
\label{cor:lin4-podobienstwo}
Niech \(V\) będzie skończenie wymiarową
przestrzenią wektorową nad \(K\),
\(T\colon V\to V\) endomorfizmem,
a \(B,B'\) dwiema uporządkowanymi bazami \(V\).
Dla \(A=[T]_{B\leftarrow B}\) i
\(P=P_{B\leftarrow B'}\) zachodzi
\[
 [T]_{B'\leftarrow B'}=P^{-1}AP.
\]
Dwie macierze kwadratowe \(A,A'\)
tego samego rozmiaru nazywamy
\emph{podobnymi}, gdy istnieje
odwracalna macierz \(P\) taka, że
\(A'=P^{-1}AP\).
\end{wniosek}
\begin{proof}
W Twierdzeniu~\ref{thm:lin4-zmiana-macierzy}
wybieramy \(W=V\), \(C=B\), \(C'=B'\).
Wtedy obie macierze przejścia są równe \(P\),
co daje wzór z tezy.
\end{proof}

\begin{przyklad}[Zmiana macierzy endomorfizmu]
\label{prz:lin4-podobienstwo}
W \(\mathbb R^2\) rozważmy
\(T(x,y)=(2x+y,y)\).
Dla bazy standardowej \(B\) mamy
\[
 A=[T]_{B\leftarrow B}
 =\begin{pmatrix}2&1\\0&1\end{pmatrix}.
\]
Niech \(B'=((1,1),(0,1))\), jak w
Przykładzie~\ref{prz:lin4-zmiana-wspolrzednych}.
Wówczas
\[
 P=\begin{pmatrix}1&0\\1&1\end{pmatrix},
 \qquad
 P^{-1}AP
 =\begin{pmatrix}1&0\\-1&1\end{pmatrix}
  \begin{pmatrix}2&1\\0&1\end{pmatrix}
  \begin{pmatrix}1&0\\1&1\end{pmatrix}
 =\begin{pmatrix}3&1\\-2&0\end{pmatrix}.
\]
Możemy to sprawdzić bezpośrednio:
\(T(1,1)=(3,1)=3(1,1)-2(0,1)\),
a \(T(0,1)=(1,1)=1(1,1)+0(0,1)\).
Nowa macierz ma więc kolumny
\((3,-2)^{\mathsf T}\) i \((1,0)^{\mathsf T}\).
\end{przyklad}

\begin{uwaga}[Co jest niezmiennicze przy zmianie bazy?]
\label{rem:lin4-niezmienniki}
Macierz odwzorowania zależy od wyboru
uporządkowanych baz, ale samo odwzorowanie
nie ulega zmianie. Jądro i obraz mają
niezależnie od baz te same wymiary:
przejścia \(P,Q\) są odwracalne,
zatem \(x\mapsto Px\) daje bijekcję
między rozwiązaniami \(Q^{-1}APx=0\)
a rozwiązaniami \(Ay=0\).
Podobnie obraz nowej macierzy
jest obrazem starej przez izomorfizm \(Q^{-1}\).
W szczególności rząd nie zmienia się
przy zmianie baz.
\end{uwaga}

\section{Transpozycja macierzy}
\label{sec:lin4-transpozycja}

\begin{definicja}[Macierz transponowana]
\label{def:lin4-transpozycja}
Niech \(m,n\in\mathbb N_0\).
Dla \(A=(a_{ij})\in\operatorname{Mat}_{m\times n}(K)\)
przez \emph{macierz transponowaną} rozumiemy
\[
 A^{\mathsf T}:=(a_{ji})_{1\leq j\leq n,\ 1\leq i\leq m}
 \in\operatorname{Mat}_{n\times m}(K).
\]
Transpozycja zamienia wiersze z kolumnami;
nie wykonuje sprzężenia zespolonego.
\end{definicja}

\begin{stwierdzenie}[Działania i transpozycja]
\label{prop:lin4-transpozycja-dzialania}
Niech \(m,n,p\in\mathbb N_0\),
\(A,B\in\operatorname{Mat}_{m\times n}(K)\)
oraz \(c\in K\). Wtedy
\[
 (A+B)^{\mathsf T}=A^{\mathsf T}+B^{\mathsf T},
 \quad(cA)^{\mathsf T}=cA^{\mathsf T},
 \quad(A^{\mathsf T})^{\mathsf T}=A.
\]
Ponadto dla \(A\in\operatorname{Mat}_{m\times n}(K)\)
i \(D\in\operatorname{Mat}_{n\times p}(K)\)
zachodzi
\[
 (AD)^{\mathsf T}=D^{\mathsf T}A^{\mathsf T}.
\]
\end{stwierdzenie}
\begin{proof}
Pierwsze trzy równości wynikają
z porównania odpowiednich wyrazów:
na przykład wyraz \((j,i)\) macierzy
\((A+B)^{\mathsf T}\) wynosi
\(a_{ij}+b_{ij}\), tak samo jak
w macierzy \(A^{\mathsf T}+B^{\mathsf T}\).

Dla ostatniego wzoru wyraz \((k,i)\)
macierzy \((AD)^{\mathsf T}\) jest równy
\[
 (AD)_{ik}=\sum_{j=1}^{n}a_{ij}d_{jk}.
\]
W macierzy \(D^{\mathsf T}A^{\mathsf T}\)
wyraz o tych samych indeksach wynosi
\(\sum_{j=1}^{n}(D^{\mathsf T})_{kj}
(A^{\mathsf T})_{ji}
=\sum_{j=1}^{n}d_{jk}a_{ij}\).
Obie wartości są równe, bo mnożenie
w ciele \(K\) jest przemienne.
Także przy pustych zbiorach indeksów
sumy i macierze są zgodne.
\end{proof}

\begin{przyklad}[Transpozycja iloczynu]
\label{prz:lin4-transpozycja}
Niech
\[
 A=\begin{pmatrix}1&2\\0&-1\end{pmatrix},
 \qquad
 D=\begin{pmatrix}2\\3\end{pmatrix}
\]
nad \(\mathbb R\).
Mamy
\(AD=(8,-3)^{\mathsf T}\), więc
\((AD)^{\mathsf T}=(8,-3)\).
Z drugiej strony
\[
 D^{\mathsf T}A^{\mathsf T}
 =\begin{pmatrix}2&3\end{pmatrix}
  \begin{pmatrix}1&0\\2&-1\end{pmatrix}
 =\begin{pmatrix}8&-3\end{pmatrix}.
\]
Przykład potwierdza odwrócenie
kolejności czynników w transpozycji iloczynu.
\end{przyklad}

\section{Układ równań liniowych jako równanie macierzowe}
\label{sec:lin4-uklady}

\begin{definicja}[Układ liniowy w postaci macierzowej]
\label{def:lin4-uklad}
Niech \(m,n\in\mathbb N_0\),
\(A=(a_{ij})\in\operatorname{Mat}_{m\times n}(K)\),
a \(b=(b_1,\ldots,b_m)^{\mathsf T}\in K^m\).
Równanie macierzowe \(Ax=b\), \(x\in K^n\),
oznacza układ \(m\) równań
\[
 \sum_{j=1}^{n}a_{ij}x_j=b_i,
 \qquad 1\leq i\leq m.
\]
Układ \(Ax=0_{K^m}\) nazywamy
\emph{jednorodnym}; inne układy nazywamy
\emph{niejednorodnymi}.
\end{definicja}

\begin{twierdzenie}[Warunek rozwiązalności i wszystkie rozwiązania]
\label{thm:lin4-zbior-rozwiazan}
Niech \(A\in\operatorname{Mat}_{m\times n}(K)\)
oraz \(b\in K^m\), gdzie \(m,n\in\mathbb N_0\).
Układ \(Ax=b\) ma rozwiązanie
wtedy i tylko wtedy, gdy \(b\in\operatorname{im}A\).
Jeżeli \(x_0\in K^n\) jest jednym rozwiązaniem,
to zbiór wszystkich rozwiązań jest równy
\[
 x_0+\ker A=\{x_0+z:z\in\ker A\}.
\]
W szczególności układ jednorodny
ma zawsze rozwiązanie zerowe.
\end{twierdzenie}
\begin{proof}
Definicja obrazu mówi, że
\(b\in\operatorname{im}A\) dokładnie wtedy,
gdy istnieje \(x_0\in K^n\) z \(Ax_0=b\).
To dowodzi pierwszej równoważności.

Załóżmy teraz \(Ax_0=b\).
Dla \(x\in K^n\) mamy
\[
 Ax=b
 \iff Ax=Ax_0
 \iff A(x-x_0)=0_{K^m}
 \iff x-x_0\in\ker A.
\]
Wykorzystaliśmy liniowość \(L_A\).
Ostatni warunek jest równoważny
istnieniu \(z\in\ker A\)
takiego, że \(x=x_0+z\).
Daje to dokładnie wskazany zbiór rozwiązań.
Przy \(b=0_{K^m}\) można wziąć \(x_0=0_{K^n}\).
\end{proof}

\begin{przyklad}[Jawne rozwiązanie układu z parametrem]
\label{prz:lin4-rozwiazania-ukladu}
Niech
\[
 A=\begin{pmatrix}1&1&0\\0&1&1\end{pmatrix},
 \qquad b=\begin{pmatrix}\alpha\\\beta\end{pmatrix},
 \quad \alpha,\beta\in\mathbb R.
\]
Układ \(Ax=b\) ma postać
\(x_1+x_2=\alpha\), \(x_2+x_3=\beta\).
Jeżeli oznaczymy \(x_2=t\in\mathbb R\),
otrzymamy
\[
 \begin{pmatrix}x_1\\x_2\\x_3\end{pmatrix}
 =\begin{pmatrix}\alpha\\0\\\beta\end{pmatrix}
 +t\begin{pmatrix}-1\\1\\-1\end{pmatrix}.
\]
Każda wartość \(t\) daje rozwiązanie,
a każde rozwiązanie ma tę postać.
Pierwszy składnik jest rozwiązaniem
szczególnym, drugi przebiega \(\ker A\)
z Przykładu~\ref{prz:lin4-rzad-prostokatnej}.
Dla każdego \(b\) układ jest rozwiązalny,
ponieważ \(\operatorname{im}A=\mathbb R^2\).
\end{przyklad}

\section{Podsumowanie}
\label{sec:lin4-podsumowanie}

Macierz odwzorowania liniowego ma w kolumnach
współrzędne obrazów wektorów bazy dziedziny.
Wzór \([T(v)]_C=[T]_{C\leftarrow B}[v]_B\)
pozwala zastąpić działanie odwzorowania
mnożeniem macierzy przez kolumnę.
Mnożenie macierzy opisuje składanie odwzorowań,
a odwracalność macierzy jest równoważna
odwracalności odpowiedniego operatora.
Przy zmianie baz macierz \(A\) zmienia się
zgodnie ze wzorem \(Q^{-1}AP\);
dla endomorfizmu jest to podobieństwo
\(P^{-1}AP\).
Rząd macierzy jest wymiarem przestrzeni
rozpiętej przez jej kolumny, a rozwiązania
układu \(Ax=b\), gdy istnieją, tworzą warstwę
\(x_0+\ker A\).
